

In spite of its philosophical intent, the media infrastructure presented here is all based on off-the-shelf information technology which ultimately comes from mines.  We must state clearly, however, that our desire is to build a media system which requires no further input of mined materials to function.  That is to say, all the media hardware we use should be based on waste streams from the existing system.  This means it can be built on discarded and broken hardware, and then as that hardware continues to disintegrate over time, that it has the means to continue to repair and evolve it so that the constituent atoms and molecules and deep structures form the basis of a stable technology which can exist in equilibrium with the local ecosystem. To build this new type of media, we need to first be clear on what the goals are, then sketch out how this might be done, create a path forward for developing this, and then take action to set in motion the series of events which will cause this system to come into being. 

When media does not have ``users'' and ``data'' as in consumer media, but is simply a network for hosting self-replicating documents, this radically alters some fundamental assumptions about how media works.  Also, when we use the waste streams of the existing wasteful system, the basic mathematics of the situation is altered by the sheer astronomical scale at which waste media hardware is currently produced. If we combine the designed-to-fail phones with the designed-to-fail laptops, and tv's and so on, and then combine on top of that all the useless junk in the so-called ``internet of things'' which is currently being accelerated, the number of physical objects we can incorporate into our new media ecosystem is absolutely staggering.  It is difficult to calculate the numbers and it is a moving target anyway, but it seems clear that we can easily build a system with at least a 10 to 1 ratio of recently-discarded media hardware to people, for every single person on this planet.  That is to say, we can conservatively estimate that we have available a reservoir of 100 billion mostly-intact media devices to distribute among a projected population of 10 billion people.  Less conservatively, and taking some of the marketing noise from Silicon Valley at face value, we might be looking at more like 100 or 1000 media devices per person, especially if we find ways to dig out a large number of more trashed older systems.

In a world with many media devices per person, and no private ``user data'' or property, we can build systems of ubiquitous community media rather than personal media devices(the primary function of which is to maximize waste, and to surveil and control people).  What does this look like?  To get a feel for how this could look we recall the discussion earlier in this book about the Street Network, and the power of physical places.  Places with significance, such as cross roads, major cultural centers, community centers, parks, fords, bridges, tunnels, libraries and so on can all host ubiquitous organic media.  

If I want a document in a place other than where I was reading and editing it before, I simply replicate that document from one network node to the next.  Documents can flow from one place to another, along with people.  But the physical infrastructure of the media does \emph{not} need to be moved from one place to another.  This is a very radical departure from existing media systems, and it alters our design constraints on hardware dramatically.  Perhaps the two most powerful effects of this difference on design is on portability and redundancy.  If we are using streams of trash to build things, and are accumulating this material in central physical network nodes, we can have a very high level of hardware redundancy.  In mine based systems, any redundancy means more money and more mining and is avoided at all cost, with the smallest amount of material in the smallest space. When devices are all attached to individual persons, they must be as small as possible, and less material is also always sought out.  But when we are absorbing material which was a liability to society, the more of it we absorb the better.  So we can build whole structures of stitched together screens for purely aesthetic purposes which never would have made sense in the old system.  Also, we can build systems that just use a single component from something and let the rest of it sit there waiting to be absorbed into some future project, but unused for the time being.  

Removing portability as a design goal makes a lot of design tasks much easier. Perhaps the most powerful consequence of this shift is with electrical power.  Today's media devices all either take power from a wall socket or from a battery which is charged via a wall socket.  Batteries are always pushed right up to the physical limits of light weight and high energy density.  But if we replace this with static media infrastructure, batteries can be replaced by hybrid power cells which store energy by several means, and take up as much space as the materials used happen to need.  We need a system for converting all the materials in wasted batteries today into working batteries.  This is not ``recycling'' in the existing centralized sense, with mass production mirroring the process of mine-based production.  Like everything in our new civilization, it is a craft mode of production, in which anyone anywhere in the world can use skills transmitted via the network to directly convert waste batteries into cells in a modular power pack which can be used to run all of our electrical devices.  Again, we note that there is no reason based on the fundamental physics and chemistry of the system that this is not possible or indeed even easy.  It is just not done today due to the broken system based on money, mining and property we are forced to work with in modern society. 

In addition to re-designing the whole power storage system of our media, we need to address power generation itself.  In order to eliminate mining and the control of empires, we need power generation which is based entirely on trash, and is totally local.  The proposed civilization here will use less than 1 100th of the power of existing civilization by removing wasted efforts.  \emph{Most} power today is used to do pointless activities, from driving around in circles to keeping lights on in giant buildings where people do pointless things all day.  A huge fraction of this is also based on our system of property which forces people to ignore weather, seasons, and natural geography and live places that make no sense to live based on the violence of the state and property.  If we ignore this and build a new civilization not tied to property, much of the energy we waste on huge megacities in unsustainable locations can be terminated as we abandon those cities(Phoenix Arizona should probably have a population of under 1000 people).  Micropower can be generated by converting all the motors found in broken electrical devices into power generators, and used for direct mechanical power generation, from wind, water, and heat engines run on the sun.  Photovoltaics and wind as they are used today are not of interest, as they rely on mine waste just like coal, nuclear and oil.  No mining means no mining.  And that means we will only build where we can get some power from the natural world above ground.

\begin{figure}
	\centering
	\includegraphics[width=3in]{figures/shapes/blank.png}
    \caption[geometronstation]
	{Geometron Station. Sketch of a Geometron Station in a Trash Camp, showing what ubiquitous organic media might look like.} 
\end{figure}


Another consequence of non-portable media is that we can scale the displays way up.  In the chapter of this book on Action Geometry, we describe a Trash Camp in which large modular structures are built for both shelters and industrial production infrastructure.  Trash camps like this which accumulate a large amount of material can have essentially all surfaces turn into screens as wasted screens pile up from the old system.  Some screens might even be dead, and just used as a flat structural element which blocks rain water.  But if we can build a modular system for integrating screens into our media, and also build out the optics technology to do large scale projection on walls and other large surfaces, we can have truly ubiquitous displays if we want to.  

In general, we also are always focusing on increasing modularity so that components which are repurposed from old broken technology have the maximum possible utility.  In many cases, this will mean sacrificing the \emph{majority} of the elements of a piece of technology and just making use of one component.  We can then, over time, as our technology improves, find ways to repeat this with other components.  But initially we will find that in many cases something like a whole smart phone is being converted to just a single modular device, say an accelerometer or a touch sensor.  In order to facilitate this kind of modular technology, we need to develop a hybrid interconnect technology in which we build up fabrication at the millimeter scale, constructing again by craft production components which break out the electrical contacts in a simple, well-documented way from needed components to the outside world.  In some cases we can build these interconnects on the existing standards like USB, but in many cases these standards are designed to make low quality stuff that fails quickly and we will need to spend considerable research efforts on building more robust standards.  Much more robust standards will be possible when we again use the principles that larger size and redundancy are now acceptable.  Without minimization of material use and size, contacts can be designed which are much more robust than the aggressively miniaturized connectors used by consumer technology.

When our network is primarily used to share documents, rather than real-time surveillance, manipulation and control, it also opens up some new possibilities in terms of information transport protocols.  The Street Network is going to create a network of travelers on the physical networks of the world: truckers, hitchhikers, camper dwellers, and wanderers of all sorts.  If we don't need documents moved instantly, this physical network can be used to transport physical storage media, creating a flow of information comparable to the trunk lines used to run the Internet today.  This network protocol is a social one rather than a technical one. Rather than information in packets being directed to numerical addresses of machines, we have documents being sent from one community of actual people to another.  This makes the protocol sound like mail, but it is not.  Mail is still based on the unique address of an individual person or organization.  The network we are building is not based on property or individuals, but on communities and on self-replication of information.  Documents will flow to where they are wanted, and actual humans will make choices on the fly about what that looks like.

In building a new type of media infrastructure, we must also consider the fundamental issues of what our machines are designed to do.  These are not computers.  We are building machines to edit, replicate, and read \emph{documents}.  The fundamental operation of a ``computer'' is arithmetic.  The fundamental operation of our machines is the display of symbols in the generalized sense discussed in an earlier chapter: any geometry with meaning.  If the only purpose of a machine is to display a document and interact with it, the whole hardware and software architecture can be completely different.  The toy model of a computer called a Turing Machine, which does things to numbers based on a program written in numbers is replaced by a model where a human makes the choice to engage a machine, requesting that the machine draw a symbol, which it does based on a program made up of geometric actions and symbols.  This is not a computer! It is Geometron: machines using symbols to make symbols rather than numbers to make numbers.

In computers, a program is running all the time, and it carries out an operation on each clock cycle, with well over a billion clock cycles per second, even when it does nothing.  In our machines, no permanent clock is needed.  When we engage with it by clicking, scrolling, typing or engaging some other sensor, we are carrying out a geometric action which causes a glyph stored in memory to be drawn.  This can have any kind of time steps, including uneven ones, where we simply want things to be done in a certain order, and each step takes as long as it takes to do the geometry.  In some cases we will want things to happen fast, but in general we aim to make these things as simple as possible, so that the things happening on a nanosecond time scale only relate to saving and loading large documents to and from memory. When we do nothing, the machine is also at rest. And it only comes alive and does things in direct response to our engagement.  All software which is not creating a document, editing one, or replicating one is dispensed with.  

The Geometron language can be used to create any of the documents described here, down to the hardware level.  We will have bitmaps be stored in an Image Stack, which is memory for the sole purpose of storing sequences of bitmaps.  A Geometron Virtual Machine can then be implemented in hardware by translating sequences of addresses into geometric actions.  These addresses reference a Geometron Hypercube which is a physical memory component which stores address sequences in each address.  With geometric movements as well as a whole font stored in the font section of the Hypercube, scrolls can be constructed as just another Geometron glyph.  Reading a book or article or any other kind of rich text document is then just a matter of drawing a single big Geometron glyph.  This glyph can include layout, graphics, text, and complex rich text structures in multiple languages.  Geometric constructions combined with access of the aforementioned Image Stack can be used to construct Maps on our displays.  If we can display Maps, Scrolls, and generic symbols, we can create and edit any document, making a generalized symbol machine in analog to the generalized computing machine.

Information is stored in self-replicating media as described in the Printers chapter: lithography in brass to make stamped copies in removable silicone films.  This information is both written and read out using old DVD drives for both the laser, the optics, and the positional control(at about 1 micron to a few cm scale), as well as the spinner to spin on the silicone coatings and photoresist layers on the polished brass plates.  Information is encoded using the Roctal coding scheme(octal that rocks).  In this scheme, numbers are represented in binary, with rows of three bits so that each row represents a digit in a base 8 system.  These arrays of bits are arranged in some standard repeating pattern with added structures for alignment marks.  This method encodes bytecode which can be used to create any of the documents described here: HTML files, JSON data, Scrolls, Maps, Feeds, Image Stacks(in base 64 encoding), and Geometron Hypercubes.  Hardware which converts these types of information into symbols can be created as a physical implementation of the GVM, directly translating sequences of addresses into sequences of geometric actions either on a display or on some physical machine.  The most fundamental such machine this controls of course is the writer.  The writer writes code onto the brass plate which represents documents, including the code to write the code.  This plate, after being etched, is used to replicate the information many times over using soft silicone films with imprints of the brass plate.  These are then read by other people on other Geometron machines, so that another person can read a set of documents which describe how to build this whole system from trash other people can find in their physical location.  Having replicated this media system, the next person who joins can then use their machines to build more self-replicating documents, by printing out brass plates and then replicating the documents again and again with one plate.  Furthermore, each instance of the system can be printing in clay at a human readable size, and we can carry out physical tabletop symbolic interactions by moving printed clay tokens around with other people.  As discussed before, these tokens are themselves self-replicating media, which can be used to stamp clay to stamp more clay and so on, replicating without a printer.  And, like the microscopic brass printed information, these tokens which represent elements of sets are used to replicate the media system by containing all the elements of the system we need to discuss when talking to someone face to face about replication.

Numerous challenging details are left out of the discussion in this chapter, but that doesn't matter.  What is clear now is that the described system can be built from a physics and engineering standpoint.  If we want it, we can build it.  So what is needed now is only to decide that we want it and that we know how we will exist in relation to it as it comes into being.  If we can do that, and if we can share this desire widely enough, it all of the technical problems will be solved by someone.  This is the full power of Trash Magic: we seek to replicate the desire to build everything from trash in both the minds of other people and in the physical reality of the world around us.  If you join in this mission, and you fully commit to feeling the desire for this world, you should be able to share that with someone else, and if they can do that, they can pass it on.  If we can spread this desire widely enough, the materials, skills, energy and time will simply materialize from the mass efforts of all these people.  Building a new civilization without money, without mining and without property is a lot easier than you think.  Now is the time for us to simply decide to do this, and get on with living the adventures that this new world will bring about.



